\section{An independent four-case proof of the 236-state fold}
\label{sec:four-case-fold}

Let $M$ be the 237-state reassociated table with SHA-256
\begin{center}\small\ttfamily
da0b3d1fdd74d56482b9d231119bbbaccdbcadc78e7b7df72c3a734a3a808724
\end{center}
and let $N$ be the new 236-state table in the inherited labels, with SHA-256
\begin{center}\small\ttfamily
ff059535876bc9e1e1e63174ebfd1c10a91c1a818ba0c3f8c4ef2ac51a3cfcc3.
\end{center}
Both start in state $95$ and halt in the distinguished symbol $254$.
The state map is
\[
 \phi(q)=\begin{cases}182,&q=223,\\q,&q\ne223.\end{cases}
\]
It changes neither the head position nor any tape bit.  This proof uses only
the literal rows listed below and the equality of the other rows.  It does
not depend on the generic symbolic fold checker or its inferred radius-five
invariant.

\subsection*{Two immediate predecessor guards}

Write a configuration as $(q,h,t)$, where $t:\mathbb Z\to\{0,1\}$.
The complete list of incoming rows to old state $182$ consists of
\[
 \delta_M(76,1)=(1,R,182),
\]
and the complete list of incoming rows to old state $223$ consists of
\[
 \delta_M(237,0)=(0,R,223).
\]
The initial state is neither $182$ nor $223$.  Consequently every reachable
old configuration satisfies the two implications
\[
 q=182\Longrightarrow t(h-1)=1,
 \qquad
 q=223\Longrightarrow t(h-1)=0.
 \tag{G}
\]
Indeed, the last transition into either state writes exactly the asserted
bit and then moves right.  This argument allows an arbitrary initial binary
tape, including tapes with infinite support, and an arbitrary initial head
position.

\subsection*{Four exact blocks}

The new state's two rows are
\[
 \delta_N(182,0)=(0,L,198),\qquad
 \delta_N(182,1)=(1,L,210).
\]
They move left to read the distinguishing bit from (G).  The following two
steps return to the original cell and then move right.  In the table, each
displayed state sequence includes the three successive new target states;
the initial new state is always $182$.

\begin{center}
\begin{tabular}{ccccc}
Old state & Current bit & Left bit & Three new target states & Old target\\
\hline
$182$ & $0$ & $1$ & $198,136,187$ & $187$\\
$182$ & $1$ & $1$ & $210,177,212$ & $212$\\
$223$ & $0$ & $0$ & $198,206,23$ & $23$\\
$223$ & $1$ & $0$ & $210,114,70$ & $70$
\end{tabular}
\end{center}

Every transition in these four blocks writes back the bit it reads.  The
head movements are $L,R,R$, so the final head is $h+1$ and the entire tape is
unchanged.  The corresponding old row also preserves the current bit and
moves right to the displayed old target.  Thus, for an old configuration
$c$ satisfying (G) in one of these two states,
\[
 N^3(\phi(c))=\phi(M(c)).
 \tag{E}
\]
None of the intermediate states is the halt symbol.

For each of the other $470$ old rows, the new row is literally the old row
with $\phi$ applied to its target.  This includes the redirection
$\delta_N(237,0)=(0,R,182)$.  Therefore
\[
 N(\phi(c))=\phi(M(c))
 \tag{O}
\]
at every other working state.  These are all $2\cdot237=474$ old rows.

\subsection*{Unbounded consequence}

Starting from $(95,h,t)$, induction on the number of old transitions applies
(E) or (O) at each step.  If $c_j$ is the old configuration after $j$ steps,
define $\tau(0)=0$ and
\[
 \tau(j+1)-\tau(j)=
 \begin{cases}3,&\text{the state of }c_j\text{ is }182\text{ or }223,\\
 1,&\text{otherwise.}\end{cases}
\]
Then $N^{\tau(j)}(95,h,t)=\phi(c_j)$ throughout the old run, and no new halt
occurs inside one of the blocks.  In particular, $j\le\tau(j)\le3j$.
If the old machine halts, the new machine reaches the same halt symbol with
exactly the same head and tape.  If the old run is infinite, $\tau(j)$ is
unbounded and its successive blocks cover the entire new run, which is
therefore infinite.  We obtain the universal theorem
\[
 \forall h\in\mathbb Z\;\forall t:\mathbb Z\to\{0,1\},\qquad
 \operatorname{Halts}(M,95,h,t)
 \;\longleftrightarrow\;
 \operatorname{Halts}(N,95,h,t).
\]

The accompanying \texttt{check\_fold236.py} independently
checks the two complete predecessor lists, all $470$ ordinary row equalities,
the four three-step blocks, the exact state sets, and the input byte hashes.
It imports neither the candidate search code nor the generic fold checker.
The proof above is ordinary finite-table mathematics formalizable in ZFC;
the file is not represented as a new proof-assistant theorem or as completion
of the inherited internal arithmetic consistency endpoint.
